\documentclass[11pt]{article}
\usepackage[utf8]{inputenc}
\usepackage[T1]{fontenc}
\usepackage{graphicx}
\usepackage{xcolor}
\usepackage{amsmath, amssymb}
\usepackage{url}
\usepackage{natbib}
\setlength{\parindent}{0em}
\usepackage[letterpaper, total={6.5in, 9in}]{geometry}

\newcommand*{\mybox}[2]{%
  \par\medskip\noindent
  \colorbox{#1!30}{\parbox{.98\linewidth}{#2}}%
  \par\medskip
}

\setlength{\parskip}{0.5em}

\usepackage{setspace}
\onehalfspacing

\usepackage{xr-hyper}   
\makeatletter
\newcommand{\myexternaldocument}[2][]{%
  \externaldocument[#1]{#2}%
  \IfFileExists{#2.aux}{}{%
    \typeout{No file #2.aux.}%  
  }%
}
\IfFileExists{newcommands.tex}{% --- LaTeX Packages ---
% --- LaTeX Packages ---
\usepackage{amsmath}
\usepackage{amsthm}
\usepackage{bm}            % For bold math symbols
\usepackage{graphicx}      % For including images
\usepackage{xcolor}        % Required for \highlight
\usepackage{xspace}        % Required for commands like \nref
\usepackage{tikz}
%\usetikzlibrary{arrows.meta}
%\usepackage{mathtools}
%\usepackage{filecontents}
\usepackage{booktabs}
%\usepackage{placeins}
\usepackage{hyperref}
        \hypersetup{
          colorlinks=true,
          linkcolor=blue,
          urlcolor=blue,
          citecolor=blue
        }
\usepackage{etoolbox}

\setlength{\bibsep}{5pt plus 5pt} % space BETWEEN items
% For tighter lines WITHIN each item:
\usepackage{setspace}
\AtBeginEnvironment{thebibliography}{\setstretch{1.05}}

% --- Custom Commands ---

% Acronyms and Custom Text
\newcommand{\CVIC}{\text{CVIC}}
\newcommand{\deviance}{\text{dev}}
\newcommand{\dpois}{\text{dpois}}
\newcommand{\ic}{\text{ic}}
\newcommand{\nmsp}{\text{\scriptsize NMSP}}
\newcommand{\nrsp}{\text{\scriptsize NRSP}}
\newcommand{\pmin}{p_{\text{\scriptsize min}}}
\newcommand{\pp}{\text{pp}}
\newcommand{\nref}{[\textbf{references}]}
\newcommand{\SP}{survival probability\xspace}
\newcommand{\SPs}{survival probabilities\xspace}
\newcommand{\svf}{S_{i}(\cdot)}
\newcommand{\unif}[0]{\text{Unif}}
\newcommand{\rsp}[0]{\text{RSP}}
\newcommand{\msp}[0]{\text{MSP}}
\newcommand{\PPP}[0]{\text{PPP}}
\newcommand{\PPPs}[0]{\text{PPPs}}
\newcommand{\PPC}[0]{\text{PPC}}
\newcommand{\bU}[0]{\bm{U}}
\newcommand{\pv}[0]{\mathcal{P}}

% Math Symbols (Bold)
\newcommand{\btheta}[0]{\bm{\theta}}
\newcommand{\bTheta}[0]{\bm{\Theta}}
\newcommand{\bx}[0]{\bm{x}}
\newcommand{\by}[0]{\bm{y}}
\newcommand{\bY}[0]{\bm{Y}}
\newcommand{\bbeta}{\bm{\beta}}
\newcommand{\bBeta}{\bm{\Beta}}
\newcommand{\blambda}{\bm{\lambda}}
\newcommand{\bmu}{\bm{\mu}}
\newcommand{\bphi}{\bm{\phi}}
\newcommand{\bp}{\bm{p}}
\newcommand{\bsgm}{\bm{\sigma}}
\newcommand{\hu}[0]{\pi^o_{i}}
\newcommand{\mhu}[0]{\pi^o}
% General Superscripts and Subscripts
\newcommand{\obs}[0]{^{\text{obs}}}
\newcommand{\sobs}[0]{^{\text{obs}}}
\newcommand{\post}[0]{^{\text{Post}}}
\newcommand{\cv}[0]{^{\text{CV}}}
\newcommand{\iscv}[0]{^{\text{ISCV}}}
\newcommand{\supt}[0]{^{(t)}}
\newcommand{\mci}{^{(t)}}        % Kept as it has a different name from \supt
\newcommand{\mcs}{^{(s)}}
\newcommand{\rep}{^{\text{rep}}}
\newcommand{\ppc}{_{\text{ppc}}}
\newcommand{\noi}{_{-i}}
\newcommand{\wi}{_{1:n}}
\newcommand{\IS}{^{\text{\scriptsize IS}}}

% Compound Superscripts for Z-residuals
\newcommand{\rspp}[0]{^{\text{Post-RSP}}}
\newcommand{\mspp}[0]{^{\text{Post-MSP}}}
\newcommand{\rspis}[0]{^{\text{ISCV-RSP}}}
\newcommand{\mspis}[0]{^{\text{ISCV-MSP}}}
\newcommand{\rsppost}{^{\text{Post-RSP }}}  % From first block
\newcommand{\rspiscv}{^{\text{ISCV-RSP}}}    % From first block
\newcommand{\E}{\mathbb{E}}
\newcommand{\Epost}{\mathbb{E}_{\btheta\,|\,\by}}
\newcommand{\Ecv}{\mathbb{E}_{\btheta\,|\,\by_{-i}}}
\newcommand{\highlight}[1]{\textcolor{blue}{#1}}
\newcommand{\pvalue}[0]{$p$-value\xspace}
\newcommand{\pvalues}[0]{$p$-values\xspace}
\newcommand{\norm}[1]{\left\lVert#1\right\rVert}

\newcommand{\textred}[1]{\textcolor{orange}{#1}}


% ===================================================================
%                 THEOREM AND PROOF DEFINITIONS
% ===================================================================



%old theorem style for ba

% % --- 1. Define Custom Theorem Styles ---
% % A style for theorems where the optional note is bold
% \newtheoremstyle{boldnote}
%   {}		    % Space above
%   {}		    % Space below
%   {\itshape}	% Body font
%   {}		    % Indent amount
%   {\bfseries}	% Theorem head font
%   {.}		    % Punctuation after head
%   {.5em}	    % Space after head
%   {\thmname{#1}\thmnumber{ #2}\thmnote{ [\textbf{#3}]}}


% % --- 2. Define Theorem-Like Environments ---
% % Set the default style for all environments defined from this point
% \theoremstyle{plain}

% % The main theorem environment, numbered per section. All others will follow its counter.
% \newtheorem{theorem}{Theorem}[section] 

% % Environments that share the theorem counter
% \newtheorem{lemma}[theorem]{Lemma}
% \newtheorem{corollary}[theorem]{Corollary}

% % Switch to the custom style for any environments you want to use it
% % \theoremstyle{boldnote}
% % \newtheorem{specialtheorem}[theorem]{Special Theorem} % Example


% % --- 3. Define Proof and Step Environments ---
% % For custom "Step" counter within proofs
% \newtheorem{proofpart}{Step}

% % Change the name "Proof" to "Proof" (bold)
% \renewcommand{\proofname}{\textbf{Proof}}

% % Automatically reset the 'Step' counter at the beginning of every proof
% \AtBeginEnvironment{proof}{\setcounter{proofpart}{0}}
% % ===================================================================
}{}

\makeatother


\myexternaldocument[S-]{supplement}


\makeatother



\begin{document}


\begin{center}
    {\LARGE \textbf{Response to Reviewer D}} \\[1em]
    {\large \today}
\end{center}
\vspace{2em} % Space between title and your text
\normalsize
\normalsize

We sincerely thank you for your time, valuable feedback, and general appreciation of the paper. Guided by the collective insights from you and other referees, we have substantially revised the manuscript. We have responded to your comments item by item, and your constructive suggestions have greatly improved both the methodology and the presentation. In this detailed response, we quote several updated passages for your reference (please note that some of these quoted excerpts may have been further refined in the final manuscript to ensure perfect flow and accuracy). In the revised manuscript, revised text is highlighted in \textred{orange}; for substantially rewritten sections, only the \textred{section title} is highlighted rather than the entire text. 

\mybox{gray}{This paper proposes a new general-purpose residual diagnostic framework for Bayesian models based on randomized survival probabilities. The authors construct $Z$-residuals that are asymptotically independent and standard normal under correct model specification. They establish asymptotic calibration for both posterior and cross-validated versions of the residuals. They also demonstrate the testing power of their approach through simulation studies. The paper shows how standard model diagnostic tools, such as QQ plots and residual-versus-covariate plots, can be applied using these $Z$-residuals.}

\paragraph{Response:} We sincerely thank the reviewer for this excellent and highly accurate summary of our manuscript. 


\medskip

\mybox{gray}{1. \textbf{Interpretation and intuition of the $Z$-residual:} In linear regression, residuals have a clear and intuitive interpretation as the difference between the observed outcome and the structural part of the model. This interpretation allows practitioners to directly link diagnostic patterns to model misspecification. In the Bayesian diagnostics framework considered in this paper, the interpretation of the proposed $Z$-residual is less transparent. It would be helpful if the authors could provide a more intuitive explanation of what the value of the $Z$-residual represents, in a manner analogous to classical residuals in linear models. Additionally, illustrations of applying the negative transformation $-\Phi^{-1}(\cdot)$ and an explanation of how this transformation helps diagnostic interpretation would be useful.}

\paragraph{Response:}
Thank you for this helpful suggestion. We agree that the interpretation of Z-residuals should be explained more intuitively.

We have revised Section~2.2, immediately after Eq.~(2), to include this intuitive interpretation and to clarify the role of the negative normal-score transformation, which is quoted as below:

\begin{quote}

This uniformity allows us to further transform the RSP into a standard normal residual via the inverse-survival function, $-\Phi^{-1}(\cdot)$, of the $N(0,1)$ distribution:
\begin{equation}\label{eq:z_residual_definition}
z_i^{\rsp} (y_i \mid \bm{\theta}) = -\Phi^{-1}(\rsp_i(y_i \mid \bm{\theta})).
\end{equation}
In words, this is the standard normal value that shares the same upper-tail area (RSP) as the observed $y_i$. Under the true model, these transformed quantities are independent $N(0,1)$ variables. We term this a ``\textbf{Z-residual}'' because $Z$ conventionally denotes a standard normal random variable. The negative sign preserves the traditional residual direction, ensuring that observations larger than expected yield positive residuals. Conceptually, a Z-residual reflects the distance from $y_i$ to its predictive median, rescaled for skewness and smoothed for discreteness. Consequently, a residual of $\pm 2$ conveys the exact same degree of extremeness regardless of the underlying predictive distribution's shape.

\end{quote}

\medskip

\mybox{gray}{2. \textbf{Practical guidance for choosing $Z$-residuals:} The paper introduces both posterior $Z$-residuals and ISCV $Z$-residuals, but it is not clear under what practical circumstances one should be preferred over the other. Providing guidance or recommendations---for example, in terms of sample size or type of model---would enhance the applicability of the proposed methodology.}

\paragraph{Response:}
Thank you for this helpful comment. We agree that practical guidance is needed for choosing between posterior and importance-sampling cross-validatory Z-residuals. 

We have revised Section 2.4 to add the following clarification: 
\begin{quote}
In practice, we generally recommend posterior Z-residuals as the default tool for routine diagnostics of models without highly localized latent effects.. While the cross-validatory construction is theoretically appealing because it alleviates the direct self-influence of $y_i$ on its own predictive distribution, it does not completely eliminate dependence in resultant LOO Z-residuals, as the training sets $\by_{-i}$ and $\by_{-j}$ share $n-2$ observations. Furthermore, ISCV Z-residuals are not exactly standard normal in finite samples---their variance is greater than 1 for high-leverage observations, akin to why frequentist ordinary least squares relies on studentized residuals. In our simulations, this structural variance, combined with importance-sampling approximation errors, resulted in higher Type I error rates for ISCV Z-residuals. Conversely, posterior Z-residuals are surprisingly more stable at smaller sample sizes, and their double-use bias vanishes asymptotically as they converge to standard normal (as justified in the appendix). Nevertheless, because ISCV computation is fast, it is good practice to compute both, as a substantial discrepancy between the two is itself informative—typically flagging highly influential observations or unstable importance-sampling approximations (e.g., large Pareto $k$ diagnostic values). Finally, ISCV Z-residuals remain distinctly superior for models with observation-specific latent variables \citep{Li2016,Millar2018}, as including $y_i$ during fitting exerts an overwhelming influence on its corresponding latent effect.
\end{quote}
In short, we recommend posterior Z-residuals as the default for routine checking. They are straightforward to compute, highly stable in finite samples, and free from importance-weighting errors. ISCV Z-residuals serve as a valuable sensitivity check when finite-sample self-influence is a concern—such as in highly flexible models or when specific observations heavily dictate the posterior.Because ISCV computation is fast, we advise calculating both. If they yield consistent results, the diagnostic conclusions are robust. A substantial discrepancy between the two is itself informative: it typically flags highly influential observations or unstable importance-sampling approximations (e.g., large Pareto $k$ diagnostic values). In such cases, analysts should inspect these influential data points and the reliability of the importance weights before finalizing their assessment.

\medskip
\mybox{gray}{3. In RSP, the random variable $U_i$ adds an additional layer of randomness to the diagnostic procedure. This raises questions regarding the stability and reproducibility of diagnostic results. It would be useful for the authors to discuss how the variability induced by this randomization affects diagnostic conclusions, and whether any strategies can be used to control this source of variability.}

\textbf{Response:} 
Thank you for this important comment. Because auxiliary randomization is required for discrete outcomes, addressing the resulting test variability is essential. To ensure stability and reproducibility, we resample the auxiliary uniforms for a fixed dataset and model, generating a distribution of replicated diagnostic $p$-values. 

Visually, diagnostic stability can be assessed via histograms of these replicated $p$-values. Quantitatively, we derive a distribution-free upper-bound $p$-value based on the median of these replications, yielding a principled, single-number summary that formally accounts for the auxiliary randomness, as reported in Sec. 2.8.

Furthermore, we have expanded Section 5 (Conclusion and Discussion) to provide practical guidance on two complementary strategies for obtaining stable single-number summaries:
\begin{enumerate}
    \item \textbf{Parameter-wise posterior predictive $p$-values (PPPs):} Our simulations demonstrate that parameter-wise Z-residual PPPs mitigate the notorious conservatism of traditional $\chi^2$ PPCs. Because they average over the MCMC draws, they remain highly stable despite the auxiliary $\bU$ drawn at each iteration, serving as a reliable secondary confirmation of model misfit.
    \item \textbf{Simulation-based calibration:} For exact calibration, analysts can simulate replicated datasets from the posterior predictive distribution to approximate the sampling distribution of the median $p$-value. This avoids the computational burden of rerunning MCMC for each replicate.
\end{enumerate}
Finally, we note that aggregating MCMC-wise $p$-values via the Cauchy combination rule \citep{liu2020cauchy} represents a highly promising strategy for handling auxiliary diagnostic variability, which we plan to explore in future research.

\vspace{1em}

\mybox{gray}{4. In the simulation study 3.2, is it possible to gain insights from any of the diagnostic results to guide users to choose the correct model family? In particular, beyond indicating that the working family is misspecified, can the diagnostic results provide interpretable patterns that suggest which alternative families may be more appropriate?}

\textbf{Response:} 
Thank you for this insightful question. Yes, diagnostic results can reveal interpretable patterns that guide users toward a more appropriate family, provided that graphical and numerical tools are evaluated jointly. While diagnostics pinpoint the \emph{direction} of model improvement, predictive information criteria are ultimately required to confirm the final model choice.

We have expanded Supplementary Section S1.2 to detail how specific visual clues suggest distinct alternative families:
\begin{itemize}
    \item \textbf{QQ plot departures:} Systematic tail deviations strongly indicate overarching distributional misspecification, suggesting alternatives like a Negative Binomial or a heavier-tailed error distribution.
    \item \textbf{Zero-inflation patterns:} In count models, if scatterplots reveal that residuals for zero observations cluster distinctly away from non-zero observations, it points users toward hurdle or zero-inflated architectures.
    \item \textbf{Covariate trends:} A pronounced trend in a LOWESS curve or a structural shift in grouped boxplots against a specific covariate usually flags a localized functional-form error (e.g., a missing nonlinear term) rather than a family misspecification.
\end{itemize}

Once these diagnostic clues point the analyst toward plausible alternatives, the final selection should be determined by comparing criteria such as the leave-one-out expected log predictive density (LOO \texttt{elpd}).

We also note an important caveat, as discussed in Section S1.2: severe misspecification in the response family can sometimes cascade and cause false rejections in localized functional-form checks (like AOV tests). Because of these cascading effects, a robust workflow requires synthesizing multiple diagnostics. Developing tests with differential sensitivity to distinct misspecification sources is an active area of our ongoing research. To maintain focus on the fundamental properties of Z-residuals relative to PPCs in this manuscript, we defer the detailed separation of zero-model and count-model misspecification diagnostics to a forthcoming paper.


\medskip

\mybox{gray}{5. Most of the numerical illustrations focus on Poisson count data. To further support the claimed generality of the framework, it would be valuable to include additional examples involving other types of discrete outcomes, such as ordinal data, to demonstrate that the proposed diagnostics perform across different data types.}

\textbf{Response:}
Thank you for this insightful suggestion. We entirely agree with your underlying premise: it is crucial to demonstrate that the proposed diagnostics perform effectively across a diverse range of data types. 

To achieve this goal and strengthen the manuscript, we opted to expand the empirical scope by introducing a continuous data setting (the inverse Gaussian regression in Section 3.2), rather than adding another discrete case like ordinal data. This choice serves a dual purpose. First, it addresses your call to demonstrate the framework's generality beyond count data. Second, because continuous data does not require auxiliary randomization, it allows us to perfectly isolate and highlight our core methodological contribution: the ability of the ``summarize-first, then-test'' paradigm to detect localized functional-form misspecifications (using covariate-specific plots and AOV tests) without the confounding variability of discrete data.

We believe this addition effectively broadens the empirical scope while powerfully contrasting the observation-level precision of Z-residuals with conventional omnibus PPCs. Because our framework indeed handles ordinal data seamlessly (requiring only standard model-implied survival probabilities and probability masses), we will feature ordinal regression examples prominently in our upcoming software package vignettes.

\setlength{\bibsep}{0.2em}
\bibliographystyle{agsm}
\bibliography{ref}

\end{document}